% Intended LaTeX compiler: pdflatex
\documentclass[a4paper,12pt]{article}
\usepackage[utf8]{inputenc}
\usepackage[T1]{fontenc}
\usepackage{graphicx}
\usepackage{longtable}
\usepackage{wrapfig}
\usepackage{rotating}
\usepackage[normalem]{ulem}
\usepackage{amsmath}
\usepackage{amssymb}
\usepackage{capt-of}
\usepackage{hyperref}
\usepackage{\string~/.emacs.d/common}
\author{mahmood}
\date{\today}
\title{dynamic programming}
\hypersetup{
 pdfauthor={mahmood},
 pdftitle={dynamic programming},
 pdfkeywords={},
 pdfsubject={},
 pdfcreator={Emacs 30.0.50 (Org mode 9.6)}, 
 pdflang={English}}
\begin{document}

\maketitle
\begin{definition}
dynamic programming amounts to breaking down an optimization problem into simpler sub-problems, and storing the solution to each sub-problem so that each sub-problem is only solved once. to be honest, this definition may not make total sense until you see an example of a sub-problem\\[0pt]
\end{definition}
\\[0pt]
\begin{question}
given the following problem:\\[0pt]
\uline{input}: an array of positive numbers of size \(n\)\\[0pt]
\uline{output}: a subarray (non-consecutive) whose sum is maximal that doesnt contain two consecutive elements\\[0pt]
example: input \(3,{\bf 12},6\) output 12\\[0pt]
example: input \({\bf 5},4,{\bf 9}\) output 14\\[0pt]
example: input \({\bf 13},3,5,{\bf 9}\) output 22\\[0pt]
example: input \({\bf 12},5,6,{\bf 17},9\) output 29\\[0pt]
\begin{answer}
we can stem the following 2 lemmas:\\[0pt]
\begin{lemma}
the optimal solution (i.e. the subarray whose sum is maximal) contains atleast a single element of every three consecutive elements\\[0pt]
\begin{proof}
we assume in contradiction that for the index i, such that \(1 \le i \le n-2\), the optimal solution doesnt contain A[i],A[i+1], and A[i+2], given that A[i+1]>0, appending A[i+1] to the given solution would give us an array with a bigger sum than the array we already have which results in a contradiction\\[0pt]
\end{proof}
\end{lemma}
\begin{lemma}
the optimal solution always contains one of A[n],A[n-1], but obviously not both\\[0pt]
\begin{proof}
if A[n-2] is not included in a given solution, then max(A[n-1],A[n]) is included, if A[n-2] is included in a solution, its obvious that a[n] is also included\\[0pt]
\end{proof}
\label{orgc18565e}
\end{lemma}
every optimal subarray (which there are a few) contains A[n-1] or A[n], if we find a subarray that contains A[n] and has the biggest sum of the arrays that contain A[n] and find a subarray that contains A[n-1] that has the biggest sum of those that contain A[n-1] then the one with the bigger sum of both is the solution\\[0pt]
to find them, we define the function \texttt{P1R} such that the call \texttt{P1R(A,i)} would find the solution for \texttt{P1} for the subarray \(A[1,\dots,i]\) under the assumption that A[i] is included in the solution\\[0pt]
for the array \texttt{A=[3,12,4,6]} we'd have \texttt{P1R(A,3) = A[1]+A[3] = 7}\\[0pt]
to understand how the function \texttt{P1R} works, we notice that given A[i] is part of the output so A[i-1] cant be included, and A[i-3] and A[i-2] cant both be included so:\\[0pt]
\[
A[i] + \max\left(\left\{P1R(A,i-2), P1R(A,i-3)\right\}\right)
\]\\[0pt]
so we finally write the algorithm:\\[0pt]
\includesvg{/Users/mahmooz/brain/out/lPe0Ruj}\\[0pt]

this solution works, but its is probably \\[0pt]
the arguments that \(\textsc{P1R}\) gets called with are the numbers \(1,2,\dots,n\), a total of \(n\) numbers, so the function gets called with the same arguments many times, to get over this and achieve a time complexity we use some sort of "cache", a \textbf{look up table}, which is what dynamic programming essentially proposes\\[0pt]
the cache we propose is an array \(B\), of size \(n\), such that \(B[i]\) is defined as the return value of \(\textsc{P1R}(A,i)\), the algorithm finds \(B[i]\) without calling \(\textsc{P1R}(A,i)\)\\[0pt]
\includesvg{/Users/mahmooz/brain/out/dZ4Qm4h}\\[0pt]
\end{answer}
\end{question}
\begin{question}
given the vector of integers A of size n, for each pair of indicies i<j we define \(\text{diff}(i,j) = A[j]-A[i]\), write an that finds \(\text{maxdiff} = \underset{i<j}{\max} \{\text{diff}(i,j)\}\}\)\\[0pt]
for example for the vector \(\{9,2,8,4,0,6,7\}\) the result should be \(\text{maxdiff} = 7 - 0 = 7\)\\[0pt]
\begin{answer}
we use an array \(B[1,\dots,n]\) such that \(B[i]=\max(A[i+1,\dots,n])\)\\[0pt]
\includesvg{/Users/mahmooz/brain/out/LhiSdNn}\\[0pt]
\end{answer}
\end{question}
\begin{question}
given a vector A of positive integers, find the maximum sum of a subarray whose indicies have a distance of more than 1\\[0pt]
\begin{answer}
we use the array B such that \(B[i]\) is the solution to the problem in the indicies \(1\dots i\) under the constraint that we have to pick \(A[i]\)\\[0pt]
\(B[i+1] = A[i+1] + \max\{B[i-1],B[i-2]\}\) is the recursion step\\[0pt]
\end{answer}
\end{question}
\begin{question}
given an array \(A\) of integers, we need to find a subarray with the maximum sum\\[0pt]
\begin{answer}
\(B[i]\) is the maximum sum that starts with \(A[i]\)\\[0pt]
\includesvg{/Users/mahmooz/brain/out/xD6Voyh}\\[0pt]

this algorithm runs in \(\Theta(n)\)\\[0pt]
\end{answer}
\end{question}
\begin{question}
given a square maze of the size \(n \times n\), we start in the upper left corner and have to arrive at the lower right corner by going down/right but not through a wall, in how many ways can this be done?\\[0pt]
\begin{answer}
we use an additional array B, such that B[i] is the number of ways to arrive from the point (i,j) in the maze to the end\\[0pt]
\end{answer}
\end{question}
\begin{question}
given an array \(A\) of s, propose an algorithm that finds the maximal sum of a subarray that doesnt contain 3 consecutive elements\\[0pt]
\begin{answer}
we use two additional arrays \(B_w\) and \(B_{wo}\), such that \(B_w[i]\) is the solution to the problem in the subarray \(A[1,\dots,i]\) such that the element \(A[i]\) is part of the sum, and \(B_{wo}[i]\) is the solution to the problem in the subarray \(A[1,\dots,i]\) without \(A[i]\) being part of the sum\\[0pt]
\includesvg{/Users/mahmooz/brain/out/QGBgbqI}\\[0pt]
\end{answer}
\label{orgecd6279}
\end{question}
\begin{my_example}
an algorithm that finds how many permutations there are of the set 0,1,2 that dont contain consecutive 1's or 2's (consecutive 0's are allowed)\\[0pt]
we use 3 additional arrays\\[0pt]
\begin{enumerate}
\item \(A_0[i]\) := the number of permutations of length i in which 0 is the last number\\[0pt]
\item \(A_1[i]\) := the number of permutations of length i in which 1 is the last number\\[0pt]
\item \(A_2[i]\) := the number of permutations of length i in which 2 is the last number\\[0pt]
\end{enumerate}
\includesvg{/Users/mahmooz/brain/out/dETdiZs}\\[0pt]

on a side note, we can use the recursive formula \(a_n=2a_{n-1}+a_{n-2}\) and to solve this\\[0pt]
\begin{note}
for now im too stupid to understand this one, ill come back to it later\\[0pt]
\end{note}
\end{my_example}
\section{matrix multiplication using dynamic programming}
\label{sec:org04c4883}
refer to for the mathematics\\[0pt]
given the matrices \(A_{i \times j} \times B_{j \times k} = C_{i \times k}\), the number of multiplication operations needed to compute \(C\) is \(i \cdot j \cdot k\), the number of addition operations needed to compute \(C\) is \(i \cdot (j-1) \cdot k\) which is close but a little smaller\\[0pt]
\begin{my_example}
given three multipliable matrices \({M_1}_{x_1 \times y_1},{M_2}_{x_2 \times y_2},{M_3}_{x_3 \times y_3}\), we want to find the number of multiplication operations required to compute \(M_{13}=M_1 \times M_2 \times M_3\)\\[0pt]
there are a few "timings" (order of application) in which we can apply the multiplication:\\[0pt]
\begin{enumerate}
\item the timing (1,2) which denotes the order of multiplication \((M_1 \times M_2) \times M_3\), which takes \(x_1x_2y_2+x_1x_3y_3=x_1x_3(x_2+y_3)\) operations\\[0pt]
\item the timing (2,1) which denotes the order \(M_1 \times (M_2 \times M_3)\) which takes \(x_2x_3y_3+x_1x_2y_3=x_2y_3(x_3+x_1)\) operations\\[0pt]
\end{enumerate}
\end{my_example}
we denote \(M_{ij} = M_i \times \dots \times M_j\) for \(1 \le i < j \le n\)\\[0pt]
the number of the operation between \(M_i\) and \(M_{i+1}\) would be \(i\)\\[0pt]
every permutation of the sequence \(S_{n-1}=1,2,\dots,n-1\) corresponds to a particular timing each with a different cost but the same resulting matrix\\[0pt]
given the multiplication of matrices \(M_{1n} = M_1 \times M_2 \times \dots \times M_n\), we define the cost of timing \(s\), \(Cost(s,M_{1n})\) as the number of multiplication operations needed to compute \(M_{1n}\) according to \(s\)\\[0pt]
the cost of finding \(M_{1n}\) is defined as the minimal cost of a given timing\\[0pt]
\[
Cost(M_{1n}) = \min_{s \in S_{n-1}}(Cost(s,M_{1n}))
\]\\[0pt]
we want an algorithm that finds the best timing (one with the minimal cost). a straightforward algorithm goes through each timing and finds the one with the minimal cost which would take \((n-1)!\) time because a timing is a of a sequence of \(n-1\) numbers\\[0pt]
we try a recursive solution. for each \((i,j),i \le j\), we denote by \(mc(i,j)\) the number of minimal multiplications required to find \(M_{ij}\). with this notation, the cost of computing \(M_{1n}\) is \(mc(1,n)\)\\[0pt]
the algorithm is based on the following 2 lemmas:\\[0pt]
\begin{lemma}
\[
(\forall 1 \le i \le j \le n)[mc(i,i)=0 \text{ and } mc(i,i+1)=x_ix_{i+1}y_{i+1}]
\]\\[0pt]
\end{lemma}
\begin{lemma}
\[
(\forall 1 \le i \le j \le n)[\text{if } j>i+1 \text{ then } mc(i,j)=\min_{k=1}^{j-1}(mc(i,k)+mc(k+1,j)+x_ix_{k+1}y_j)
\]\\[0pt]
\end{lemma}
\includesvg{/Users/mahmooz/brain/out/do6iLIn}\\[0pt]

if we were to follow the steps of the algorithm, we'd notice the following:\\[0pt]
to find mc(1,5) the algorithm has to find the minimum of:\\[0pt]
\begin{gather*}
  mc(1,1) + mc(2,5) + x_1x_2y_5\\
  mc(1,2) + mc(3,5) + x_1x_3y_5\\
  mc(1,3) + mc(4,5) + x_1x_4y_5\\
  mc(1,4) + mc(5,5) + x_1x_5y_5
\end{gather*}
and to find mc(2,5) we'd need to find the minimum of:\\[0pt]
\begin{gather*}
  mc(2,2) + mc(3,5) + x_2x_3y_5\\
  mc(2,3) + mc(4,5) + x_2x_4y_5\\
  mc(2,4) + mc(5,5) + x_2x_5y_5
\end{gather*}
we'd notice for example that mc(3,5) and mc(4,5) are being calculated twice, so the algorithm does some repeated work that can be reduced using dynamic programming\\[0pt]
the algorithm has a lower bound of \(2^{n-1}\) and while it is better than the straightforward solution, it is still inefficient\\[0pt]
each of the two arguments that the function \texttt{mc} takes can vary between n values, so the total number of permutations of arguments it can take is \(n^2\), so we use a \href{20230208201152-array_as_map.org}{lookup table} of size \(n^2\) and use it to store the results of executing the function \texttt{mc} with each possible configuration of the arguments\\[0pt]
the explanation is in \href{file:///Volumes/brainpartition/brain/notes/data/C6/4FF75D-6235-43D2-9CB5-601918A73D95/הרצאה 8 אלגוריתמים א חוברת.pdf}{file for lecture8} (too incurious to understand/document this), using \href{20230102173117-dynamic_programming.org}{dynamic programming}, we come up with this algorithm:\\[0pt]
\includesvg{/Users/mahmooz/brain/out/uedvPld}\\[0pt]

this algorithm runs in \(\Theta(n^3)\)
\end{document}